\begin{figure}[H]
\centering
\begin{tikzpicture}[scale=1.15]
  \coordinate (O) at (0,0);
  \coordinate (U) at (3.5,2.5);
  \coordinate (V) at (5,1);
  \coordinate (P) at (3.846,0.769);
  \draw[-{Stealth},very thick,uniblue] (O) -- (V) node[right] {$\ket v$};
  \draw[-{Stealth},very thick,unired] (O) -- (U) node[above] {$\ket u$};
  \draw[-{Stealth},very thick,unigreen] (O) -- (P)
        node[midway,below,font=\small] {$\dfrac{\braket vu}{\braket vv}\ket v$};
  \draw[dashed,thick,black] (P) -- (U)
        node[midway,left=3pt,font=\small]
        {$\ket\phi$};
  \draw (4.140,0.828) -- (4.081,1.122) -- (3.787,1.063);
  \draw pic["$\vartheta$",draw=unigray,angle radius=1.1cm,angle eccentricity=1.35]
        {angle=V--O--U};
  \node[font=\small,align=left] at (7.2,3.0)
    {$\ket\phi=\ket u-\dfrac{\braket vu}{\braket vv}\ket v$\\[4pt]
     $\braket\phi v=0$ \\[2pt] $\Rightarrow\ \braket\phi\phi\ge0$ \\[2pt]
     $\Rightarrow\ |\braket uv|^2\le\braket uu\braket vv$};
\end{tikzpicture}
\caption{Cauchy--Schwarz: $\lambda=-\braket vu/\braket vv$ hace que $\ket\phi=\ket u+\lambda\ket v$ sea el residuo ortogonal a $\ket v$; su norma no puede ser negativa.}
\label{fig:cauchy_schwarz}
\end{figure}